\documentclass[11pt]{article}
\usepackage[utf8]{inputenc}
\usepackage{amsmath,amssymb,amsthm}
\usepackage[margin=1in]{geometry}

\newtheorem{lemma}{Lemma}
\newtheorem{theorem}{Theorem}

\title{Problem 686: Powers of Two}
\author{}
\date{}

\begin{document}
\maketitle

\section{Problem Statement}

Find $p(123,678910)$, the $678910$-th smallest exponent $j$ such that the decimal expansion of $2^j$ begins with the digits \texttt{123}.

\section{Mathematical Foundation}

Let $\alpha = \log_{10} 2$.

\begin{lemma}
The number $2^j$ begins with the digits \texttt{123} if and only if
\[
  \log_{10}(1.23) \le \{j\alpha\} < \log_{10}(1.24),
\]
where $\{x\}$ denotes the fractional part of $x$.
\end{lemma}

\begin{proof}
We have the leading-digit condition exactly when there exists an integer $m$ such that
\[
  123 \cdot 10^m \le 2^j < 124 \cdot 10^m.
\]
Taking base-$10$ logarithms gives
\[
  \log_{10}(123) + m \le j\alpha < \log_{10}(124) + m.
\]
Subtracting $2+m$ throughout yields
\[
  \log_{10}(1.23) \le \{j\alpha\} < \log_{10}(1.24).
\]
The converse is immediate by reversing the same steps.
\end{proof}

\begin{theorem}
Scanning the exponents $j=1,2,3,\dots$ and counting those satisfying the interval test from the lemma yields $p(123,678910)$.
\end{theorem}

\begin{proof}
The interval test is equivalent to the defining property, so the $678910$-th hit is exactly the required exponent.
\end{proof}

\section{Algorithm}

\begin{enumerate}
  \item Compute $\alpha = \log_{10}2$, $L=\log_{10}(1.23)$, and $U=\log_{10}(1.24)$.
  \item For $j=1,2,3,\dots$, compute the fractional part $\{j\alpha\}$.
  \item Count those $j$ for which $L \le \{j\alpha\} < U$.
  \item Stop at the $678910$-th such exponent.
\end{enumerate}

\section{Complexity Analysis}

\begin{itemize}
  \item Time: $O(p(123,678910))$.
  \item Space: $O(1)$.
\end{itemize}

\section{Answer}

\[
\boxed{193060223}
\]

\end{document}
